\documentclass[11pt]{article}
\usepackage[margin=1in]{geometry}
\usepackage{amsmath, amssymb, amsthm}
\usepackage{mathtools}
\usepackage{booktabs}
\usepackage{xcolor}
\usepackage{hyperref}

% Macros matching RASC-MeetingNotes.tex
\newcommand{\bx}{\mathbf{x}}
\newcommand{\by}{\mathbf{y}}
\newcommand{\proj}{\operatorname{proj}}
\newcommand{\prox}{\operatorname{prox}}
\newcommand{\TV}{\mathrm{TV}}
\newcommand{\amin}{\mathop{\mathrm{argmin}}}
\newcommand{\mini}{\mathop{\mathrm{minimize}}}
\newcommand{\stn}{\text{s.t.}}
\newcommand{\R}{\mathbb{R}}
\newcommand{\N}{\mathcal{N}}

\theoremstyle{plain}
\newtheorem{proposition}{Proposition}
\newtheorem{lemma}{Lemma}
\theoremstyle{remark}
\newtheorem{remark}{Remark}

\begin{document}

\begin{center}
  {\Large\bfseries Proximal-Gradient Residual for the $y$-Block:\\
  Derivation and Justification of Equation~(9)}

  \medskip
  {\em A technical note accompanying \texttt{RASC-MeetingNotes.tex}}

  \smallskip
  Sven Leyffer and Dominic Yang, Argonne National Laboratory

  \smallskip
  \today
\end{center}

\tableofcontents

%%%%%%%%%%%%%%%%%%%%%%%%%%%%%%%%%%%%%%%%%%%%%%%%%%%%%%%%%%%%%%%%%%%%%%%%%%%%
\section{Motivation and Scope}

The ADMM-Filter algorithm of \texttt{RASC-MeetingNotes.tex} uses a scalar
\emph{dual infeasibility measure} $\widetilde{\omega}_\rho(x,y,\lambda)$ to
decide whether an iterate $(x,y)$ is a candidate for filter acceptance and
for termination (when $\widetilde{\omega}_0 \leq \varepsilon_{\mathrm{opt}}$).
The design requirement is that $\widetilde{\omega}_\rho = 0$ if and only if
$(x,y,\lambda)$ is a first-order stationary point of the augmented Lagrangian
$\mathcal{L}_\rho$ over the feasible set.

The natural candidate---the \emph{projected-gradient residual} from classical
smooth optimality theory---breaks down because the $y$-block
of $\mathcal{L}_\rho$ contains the non-smooth graph Total Variation
$\alpha\,\TV(y)$.
This note explains:
\begin{enumerate}
  \item[(i)] What the projected-gradient residual is and why it cannot be
             applied directly to the $y$-block (Section~\ref{sec:smooth}).
  \item[(ii)] How the proximal operator replaces the projection in the
              non-smooth case and gives the correct fixed-point
              characterisation (Section~\ref{sec:prox}).
  \item[(iii)] Why equation~(9) of \texttt{RASC-MeetingNotes.tex} is the
               correct residual for the $y$-block, with a proof that it
               vanishes exactly at stationary points (Section~\ref{sec:eq9}).
  \item[(iv)] How the proximal operator in equation~(9) is computed in
              practice via the Chambolle--Pock primal-dual algorithm
              (Section~\ref{sec:computation}).
  \item[(v)] How the two block residuals combine into
             $\widetilde{\omega}_\rho$ and the role of the step size $\tau$
             (Section~\ref{sec:omega}).
\end{enumerate}

%%%%%%%%%%%%%%%%%%%%%%%%%%%%%%%%%%%%%%%%%%%%%%%%%%%%%%%%%%%%%%%%%%%%%%%%%%%%
\section{Problem Setup}

Recall the augmented Lagrangian (equation~(2) of \texttt{RASC-MeetingNotes.tex}):
\begin{equation}\label{eq:aug-lag}
  \mathcal{L}_\rho(x,y,\lambda)
  = F(x) + \alpha\,\TV(y)
    - \lambda^\top(x - y)
    + \tfrac{\rho}{2}\|x - y\|_2^2,
\end{equation}
where $F:\R^n\to\R$ is the smooth structural compliance, $\TV(y) =
\sum_{(u,v)\in E}|y_u - y_v|$ is the graph total variation over element
neighbours (non-smooth and non-differentiable), $\lambda\in\R^n$ is the ADMM
dual variable, and $\rho > 0$ is the penalty parameter.  The feasible set is
\[
  \mathcal{X} \times \mathcal{Y}
  := \bigl\{(x,y): x,y\in[\epsilon,1]^n,\;\mathbf{1}^\top x \leq B,\;
    \mathbf{1}^\top y \leq B\bigr\}.
\]

The dual infeasibility $\omega_\rho$ is meant to measure how far $(x,y)$
is from satisfying the first-order stationarity conditions of
\[
  \mini_{(x,y)\in\mathcal{X}\times\mathcal{Y}}\;
  \mathcal{L}_\rho(x,y,\lambda)
\]
for a fixed $\lambda$.  Since the blocks $\mathcal{X}$ and $\mathcal{Y}$
decouple, it suffices to measure stationarity of each block separately.
The $x$-block contains only smooth terms and can be handled by a classical
projected-gradient residual (Section~\ref{sec:smooth}).  The $y$-block is
non-smooth and requires the proximal extension (Section~\ref{sec:prox}).

%%%%%%%%%%%%%%%%%%%%%%%%%%%%%%%%%%%%%%%%%%%%%%%%%%%%%%%%%%%%%%%%%%%%%%%%%%%%
\section{The Smooth Case: Projected-Gradient Residual}
\label{sec:smooth}

\subsection{Projected-gradient residual for smooth constraints}

For a \emph{smooth} function $f:\R^n\to\R$ and a closed convex set $C\subset\R^n$,
the first-order optimality condition for $\mini_{z\in C} f(z)$ is
\begin{equation}\label{eq:kkt-smooth}
  -\nabla f(z^*) \in \mathcal{N}_C(z^*),
\end{equation}
where $\mathcal{N}_C(z^*)$ is the normal cone to $C$ at $z^*$.  A classical
equivalent characterisation is the \emph{projected-gradient fixed-point}:
for any $\tau > 0$,
\begin{equation}\label{eq:proj-fp}
  z^* \;\text{solves}\;\eqref{eq:kkt-smooth}
  \quad\iff\quad
  z^* = \proj_C\!\bigl[z^* - \tau\,\nabla f(z^*)\bigr].
\end{equation}
This follows immediately from the non-expansive projection: $z^* = \proj_C(v)$
iff $v - z^* \in \mathcal{N}_C(z^*)$.  See \cite{ParBoy14,Bertsekas99} for
standard treatments.

The \emph{projected-gradient residual} at an arbitrary point $z$ is
\begin{equation}\label{eq:proj-res}
  r_{\mathrm{pg}}(z) := \proj_C\!\bigl[z - \tau\,\nabla f(z)\bigr] - z.
\end{equation}
By \eqref{eq:proj-fp}, $r_{\mathrm{pg}}(z^*) = 0$ iff $z^*$ is optimal.
Since $r_{\mathrm{pg}}$ is Lipschitz continuous in $z$ for smooth $f$, its
squared norm $\|r_{\mathrm{pg}}(z)\|_2^2$ is a continuous, non-negative
optimality measure.

\subsection{Application to the $x$-block}

The $x$-block of $\mathcal{L}_\rho$ at fixed $(y,\lambda)$ is:
\[
  \mini_{x\in\mathcal{X}}\;
  \underbrace{F(x) - \lambda^\top x + \tfrac{\rho}{2}\|x - y\|_2^2}_{=:\,h(x)},
\]
where $h$ is smooth ($F$ is smooth, the rest is a smooth quadratic).  Its
gradient is
\[
  \nabla_x \mathcal{L}_\rho
  = \nabla F(x) - \lambda + \rho(x - y),
\]
and the $x$-block residual of equation~(8) in \texttt{RASC-MeetingNotes.tex}
is exactly \eqref{eq:proj-res} applied to $h$ over $\mathcal{X}$:
\begin{equation}\label{eq:r_x}
  r_x(x,y,\lambda)
  := \proj_{\mathcal{X}}\!\bigl[x - \tau\,\nabla_x \mathcal{L}_\rho\bigr] - x.
\end{equation}
This is well-defined and $r_x = 0$ iff $x$ satisfies the first-order
conditions for the $x$-block.

\subsection{Why the projected gradient fails for the $y$-block}

The $y$-block of $\mathcal{L}_\rho$ at fixed $(x,\lambda)$ is:
\[
  \mini_{y\in\mathcal{Y}}\;
  \underbrace{\alpha\,\TV(y)}_{\text{non-smooth}}
  + \underbrace{\lambda^\top y - \tfrac{\rho}{2}\|y - x\|_2^2 + \mathrm{const}}_{\text{smooth}}.
\]
The projected-gradient residual would require $\nabla_y \mathcal{L}_\rho$,
but $\TV(y)$ is non-differentiable whenever $y_u = y_v$ for some edge
$(u,v)\in E$---a condition that is \emph{generically satisfied} at an optimal
design (piece-wise constant density fields lie in the non-differentiable
region of TV).  Applying the projected gradient here would either require an
ad-hoc subgradient (whose choice affects the measure and breaks the
fixed-point property) or ignoring the TV term in the gradient entirely
(which would give a measure that can be zero at non-stationary points).

The correct generalisation is the \emph{proximal-gradient} extension.

%%%%%%%%%%%%%%%%%%%%%%%%%%%%%%%%%%%%%%%%%%%%%%%%%%%%%%%%%%%%%%%%%%%%%%%%%%%%
\section{Proximal Operators and the Non-Smooth Fixed-Point}
\label{sec:prox}

\subsection{Definition of the proximal operator}

For a proper, closed, convex function $\phi:\R^n\to(-\infty,+\infty]$, the
\emph{proximal operator} with step size $\tau > 0$ is \cite{ParBoy14}:
\begin{equation}\label{eq:prox-def}
  \prox_{\tau\phi}(v)
  := \amin_{z\in\R^n}\;
    \tau\,\phi(z) + \tfrac{1}{2}\|z - v\|_2^2.
\end{equation}
This is a well-defined single-valued map (the objective is strictly convex in
$z$).  Two key special cases:
\begin{itemize}
  \item When $\phi = \iota_C$ (the indicator of a closed convex set $C$),
    $\prox_{\tau\phi}(v) = \proj_C(v)$.
  \item When $\phi = 0$, $\prox_{\tau\phi}(v) = v$.
\end{itemize}
Hence the proximal operator generalises both the Euclidean projection and the
identity map \cite{ParBoy14,RyuBoy16}.

\subsection{The proximal fixed-point characterisation}

The following proposition is the non-smooth analogue of \eqref{eq:proj-fp};
it is the basis for the proximal gradient method \cite{BecTeb09} and is
proved in many standard references \cite{ParBoy14,Bec17}.

\begin{proposition}[Proximal fixed-point]\label{prop:prox-fp}
  Let $g:\R^n\to\R$ be convex and $C^1$, and let $h:\R^n\to(-\infty,+\infty]$
  be proper, closed, and convex.  Define $\phi := h + \iota_C$ for a closed
  convex set $C$.  For any $\tau > 0$,
  \begin{equation}\label{eq:prox-fp}
    z^* \in \amin_{z\in C}\bigl[g(z) + h(z)\bigr]
    \quad\iff\quad
    z^* = \prox_{\tau(h + \iota_C)}\!\bigl[z^* - \tau\,\nabla g(z^*)\bigr].
  \end{equation}
\end{proposition}

\begin{proof}
  By the definition of the proximal operator, $z^* = \prox_{\tau\phi}(v)$ iff
  $v - z^* \in \tau\,\partial\phi(z^*)$ (the subdifferential characterisation
  of the proximal map; see \cite{ParBoy14} Proposition~2.2).  Setting
  $v := z^* - \tau\,\nabla g(z^*)$:
  \begin{align*}
    &z^* = \prox_{\tau(h+\iota_C)}\!\bigl[z^* - \tau\,\nabla g(z^*)\bigr]\\
    &\quad\iff\;
    \bigl(z^* - \tau\nabla g(z^*)\bigr) - z^* \in \tau\,\partial(h + \iota_C)(z^*)\\
    &\quad\iff\;
    -\nabla g(z^*) \in \partial h(z^*) + \mathcal{N}_C(z^*),
  \end{align*}
  which is exactly the first-order optimality condition for
  $\mini_{z\in C}[g(z) + h(z)]$:
  \[
    0 \in \nabla g(z^*) + \partial h(z^*) + \mathcal{N}_C(z^*). \qedhere
  \]
\end{proof}

\begin{remark}
  The condition \eqref{eq:prox-fp} holds for \emph{any} $\tau > 0$; the
  choice of $\tau$ affects the magnitude of the residual at non-optimal points
  but not the location of the zeros.  In \texttt{optimality.py} the default
  is $\tau = 1$.
\end{remark}

%%%%%%%%%%%%%%%%%%%%%%%%%%%%%%%%%%%%%%%%%%%%%%%%%%%%%%%%%%%%%%%%%%%%%%%%%%%%
\section{Equation~(9): The $y$-Block Proximal-Gradient Residual}
\label{sec:eq9}

\subsection{Splitting the $y$-block into smooth and non-smooth parts}

Write the $y$-block objective as $g(y) + h(y)$ where
\begin{align*}
  g(y) &:= -\lambda^\top y + \tfrac{\rho}{2}\|y - x\|_2^2
           \quad\text{(smooth, strongly convex in $y$)},\\
  h(y) &:= \alpha\,\TV(y)
           \quad\text{(non-smooth, convex)}.
\end{align*}
The smooth gradient is (equation~(7) of \texttt{RASC-MeetingNotes.tex}):
\begin{equation}\label{eq:smooth-grad}
  \nabla g(y) = \nabla_y \mathcal{L}^{\mathrm{sm}}_\rho(x,y,\lambda)
  = \lambda + \rho(y - x).
\end{equation}

\subsection{The residual and its zero-set}

Applying Proposition~\ref{prop:prox-fp} with $C = \mathcal{Y}$,
$h = \alpha\,\TV$, and the smooth gradient \eqref{eq:smooth-grad}, equation~(9)
of \texttt{RASC-MeetingNotes.tex} defines the \emph{$y$-block
proximal-gradient residual}:
\begin{equation}\label{eq:r_y}
  r_y(x,y,\lambda)
  := \prox_{\tau(\alpha\,\TV + \iota_{\mathcal{Y}})}\!\bigl[
       y - \tau\,\nabla_y \mathcal{L}^{\mathrm{sm}}_\rho(x,y,\lambda)
     \bigr] - y.
\end{equation}

\begin{proposition}\label{prop:ry-zero}
  $r_y(x,y,\lambda) = 0$ if and only if $y$ satisfies the first-order
  optimality conditions for $\mini_{y\in\mathcal{Y}}
  \bigl[\alpha\,\TV(y) + g(y)\bigr]$ at fixed $(x,\lambda)$,
  i.e.\ $y$ minimises the $y$-block of $\mathcal{L}_\rho$.
\end{proposition}
\begin{proof}
  Immediate from Proposition~\ref{prop:prox-fp} with the identifications above.
\end{proof}

In other words, $r_y$ is an exact optimality certificate for the $y$-block:
it vanishes precisely at the points where $y$ would be the output of a perfect
(exact) Chambolle--Pock solve.  At non-optimal iterates it provides a
continuous, non-negative quantity that the filter algorithm uses to gauge
distance from dual feasibility.

\subsection{Connection to the classical projected-gradient residual}

When $\tau = 1$ and the non-smooth term is absent ($\alpha = 0$, so
$h \equiv 0$), the prox in \eqref{eq:r_y} reduces to $\proj_{\mathcal{Y}}$,
and \eqref{eq:r_y} collapses to the classical projected-gradient residual
$\proj_{\mathcal{Y}}[y - \nabla g(y)] - y$.  Thus equation~(9) is a strict
generalisation of the classical form.

The same argument applies to the $x$-block: since $F$ is smooth, the
projected-gradient residual \eqref{eq:r_x} is recovered from
Proposition~\ref{prop:prox-fp} with $h \equiv 0$ and $C = \mathcal{X}$.
Choosing $\tau = 1$ for both blocks then recovers the structure of
equation~(6) in \texttt{RASC-MeetingNotes.tex} (the original smooth form
of $\omega_\rho$).

%%%%%%%%%%%%%%%%%%%%%%%%%%%%%%%%%%%%%%%%%%%%%%%%%%%%%%%%%%%%%%%%%%%%%%%%%%%%
\section{Computing the Proximal Operator via Chambolle--Pock}
\label{sec:computation}

\subsection{The proximal subproblem}

By definition \eqref{eq:prox-def}, evaluating $r_y$ at a point $y$ requires
solving
\begin{equation}\label{eq:prox-problem}
  z^* = \prox_{\tau(\alpha\,\TV + \iota_{\mathcal{Y}})}(v),
  \qquad v := y - \tau\,\nabla_y \mathcal{L}^{\mathrm{sm}}_\rho,
\end{equation}
which expands to the constrained TV denoising problem:
\begin{equation}\label{eq:tv-denoise}
  \mini_{z\in\mathcal{Y}}\;
  \tau\alpha\,\TV(z) + \tfrac{1}{2}\|z - v\|_2^2.
\end{equation}

\subsection{Matching the Chambolle--Pock interface}

The module \texttt{graph\_tv.py} solves problems of the form
\begin{equation}\label{eq:cp-interface}
  \mini_{z\in\mathcal{Y}}\;
  \alpha'\,\TV(z) + \sum_v \bigl(a_v\,z_v^2 + b_v\,z_v\bigr),
\end{equation}
using the primal-dual algorithm of \cite{ChaPoc11}.  Matching \eqref{eq:tv-denoise}
with \eqref{eq:cp-interface} via $(1/2)\|z - v\|^2 = (1/2)z^\top z - v^\top z + \mathrm{const}$
gives:
\[
  \alpha' = \tau\alpha,
  \qquad
  a_v = \tfrac{1}{2}\;\forall v,
  \qquad
  b_v = -v_v\;\forall v,
\]
which is exactly the call in \texttt{optimality.py} (lines 111--122):
\begin{verbatim}
  v = y - tau * grad_y_L_smooth
  y_step, _ = chambolle_pock_graph_tv(
      n_vertices=n, edges=edges,
      a=0.5 * np.ones(n),
      b=-v,              # b_v = -v_v
      budget=budget,
      alpha=tau * alpha, # alpha' = tau * alpha
      x_lo=x_lo, ...
  )
  r_y = y_step - y
\end{verbatim}

\subsection{The Chambolle--Pock algorithm for TV subproblems}

The Chambolle--Pock algorithm \cite{ChaPoc11} is a first-order primal-dual
method for problems of the form $\min_z G(z) + H(Kz)$, where $K$ is a linear
operator and $G,H$ are proper closed convex functions.  For the TV denoising
problem \eqref{eq:tv-denoise}, set $K$ to be the (signed) edge-incidence
matrix of the mesh graph (so that $(Kz)_{(u,v)} = z_u - z_v$), $G(z) =
(1/2)\|z-v\|^2 + \iota_{\mathcal{Y}}(z)$, and $H(w) = \tau\alpha\|w\|_1$.
Then the algorithm alternates:
\begin{enumerate}
  \item \emph{Primal step}: $z^{(t+1)} = \prox_{\sigma G}(z^{(t)} - \sigma K^\top w^{(t)})$,
    which is a box-and-budget projection followed by a quadratic solve
    (closed form via the bisection approach of \texttt{projection.py}).
  \item \emph{Dual step}: $w^{(t+1)} = \prox_{\sigma H^*}(w^{(t)} + \sigma K z^{(t+1)})$,
    where $H^*$ denotes the \emph{convex conjugate} (Fenchel conjugate) of $H$.
    Since $H(w) = \tau\alpha\|w\|_1$, its conjugate is
    $H^*(s) = \sup_w\{s^\top w - \tau\alpha\|w\|_1\} = \iota_{\|s\|_\infty \leq \tau\alpha}(s)$,
    the indicator of the $\ell^\infty$ ball of radius $\tau\alpha$.
    Hence $\prox_{\sigma H^*} = \proj_{\|s\|_\infty \leq \tau\alpha}$,
    which clips the edge flows componentwise to $[-\tau\alpha, \tau\alpha]$.
\end{enumerate}
Step sizes satisfy $\sigma\tau \leq 1/\|K\|^2$ to guarantee convergence
\cite{ChaPoc11}; in \texttt{graph\_tv.py} they are chosen as $\sigma = \tau_{\mathrm{cp}}
= 1/\sqrt{2d_{\max} + n}$ where $d_{\max}$ is the maximum degree of the mesh
graph, since $\|K\|^2 \leq 2d_{\max}$ for a graph incidence matrix.

\begin{remark}[Warm starting]
  The Chambolle--Pock call in \texttt{optimality.py} is warm-started at
  $z^{(0)} = y$, since $y$ is typically close to the solution of the prox
  subproblem.  This reduces the inner iteration count substantially at later
  outer iterations when $y \approx y^*$.
\end{remark}

%%%%%%%%%%%%%%%%%%%%%%%%%%%%%%%%%%%%%%%%%%%%%%%%%%%%%%%%%%%%%%%%%%%%%%%%%%%%
\section{The Combined Measure $\widetilde{\omega}_\rho$ and the Role of $\tau$}
\label{sec:omega}

\subsection{Definition}

Combining the $x$-block residual \eqref{eq:r_x} and the $y$-block residual
\eqref{eq:r_y}, the proximal-gradient dual infeasibility (equation~(10) of
\texttt{RASC-MeetingNotes.tex}) is:
\begin{equation}\label{eq:omega-prox}
  \widetilde{\omega}_\rho(x,y,\lambda)
  := \|r_x(x,y,\lambda)\|_2^2 + \|r_y(x,y,\lambda)\|_2^2.
\end{equation}
By Propositions~\ref{prop:ry-zero} and the analogous statement for the
$x$-block, $\widetilde{\omega}_\rho = 0$ iff $(x,y)$ is simultaneously
stationary for both blocks of $\mathcal{L}_\rho$ at fixed $\lambda$, i.e.\
$(x,y)$ is a saddle point of $\mathcal{L}_\rho$.

\subsection{The identity relating $\widetilde{\omega}_\rho$ and $\widetilde{\omega}_0$}

Let $\lambda^{k+1} = \lambda^k + \rho(y - x)$ be the standard ADMM multiplier
update.  A direct computation shows that the smooth gradient at the new
multiplier equals the penalised gradient at the old multiplier:
\begin{align*}
  \nabla_x \mathcal{L}^{\mathrm{sm}}_0\big|_{\lambda^{k+1}}
    &= \nabla F(x) - \lambda^{k+1}
     = \nabla F(x) - \lambda^k - \rho(y-x)
     = \nabla_x \mathcal{L}^{\mathrm{sm}}_\rho\big|_{\lambda^k},\\
  \nabla_y \mathcal{L}^{\mathrm{sm}}_0\big|_{\lambda^{k+1}}
    &= \lambda^{k+1}
     = \lambda^k + \rho(y-x)
     = \nabla_y \mathcal{L}^{\mathrm{sm}}_\rho\big|_{\lambda^k}.
\end{align*}
Both blocks see the same gradient, so the residuals are identical:
\begin{equation}\label{eq:identity}
  \widetilde{\omega}_\rho(x,y,\lambda^k)
  = \widetilde{\omega}_0(x,y,\lambda^{k+1}).
\end{equation}
This is the identity cited in \S\ref{subsec:omega-zero} of
\texttt{RASC-MeetingNotes.tex} and justifies evaluating the post-multiplier-update
convergence test with $\rho = 0$: the $\varepsilon_{\mathrm{opt}}$-test
targets $\widetilde{\omega}_0$, the unpenalised first-order error, which by
\eqref{eq:identity} is available for free from the last inner-loop
$\widetilde{\omega}_\rho$ value.

\subsection{Role of the step size $\tau$}

The step size $\tau > 0$ in equations~(8)--(9) and \eqref{eq:omega-prox} is
a free parameter.  Its choice does not affect \emph{where} the residuals
vanish (the optimal set is fixed by the problem, not by $\tau$), but it
scales the residuals at non-optimal points.  Two considerations:
\begin{itemize}
  \item With $\tau = 1$, equation~(8) recovers the classical projected-gradient
    residual on the $x$-block and equation~(9) matches the structure of the
    original $\omega_\rho$ (equation~(6)) for the smooth case.  This is the
    default in \texttt{optimality.py}.
  \item A penalty-adaptive choice such as $\tau = 1/\rho$ or $\tau = 1/(\rho +
    L_F)$ (where $L_F$ is the Lipschitz constant of $\nabla F$) would
    normalise the magnitude of $\widetilde{\omega}_\rho$ across the
    $\rho$-doublings in the restoration phase and might improve the filter
    acceptance constants.  This is noted as a known issue in
    \texttt{STATUS.md} (\S ``Known issues'', $\tau = 1.0$ item).
\end{itemize}

%%%%%%%%%%%%%%%%%%%%%%%%%%%%%%%%%%%%%%%%%%%%%%%%%%%%%%%%%%%%%%%%%%%%%%%%%%%%
\section{Summary}

\begin{center}
\begin{tabular}{@{}lll@{}}
\toprule
\textbf{Block} & \textbf{Smoothness} & \textbf{Residual at point $z$}\\
\midrule
$x$-block (eq.~(8)) & $F$ smooth &
  $r_x = \proj_{\mathcal{X}}\bigl[x - \tau\nabla_x \mathcal{L}_\rho\bigr] - x$ \\[4pt]
$y$-block (eq.~(9)) & $\alpha\,\TV$ non-smooth &
  $r_y = \prox_{\tau(\alpha\TV+\iota_{\mathcal{Y}})}\bigl[y-\tau\nabla_y \mathcal{L}^{\mathrm{sm}}_\rho\bigr] - y$ \\
\bottomrule
\end{tabular}
\end{center}

\medskip\noindent
Equation~(9) is the correct non-smooth generalisation because
(a) its zero-set coincides exactly with the first-order stationary points of
the $y$-block (Proposition~\ref{prop:ry-zero}), (b) it reduces to the
projected-gradient residual when the TV term is absent, (c) the proximal
operator it requires is already solved by the existing Chambolle--Pock
routine with a simple reparameterisation ($a_v = 1/2$, $b_v = -v_v$,
$\alpha' = \tau\alpha$), and (d) the identity~\eqref{eq:identity} provides
the post-multiplier-update convergence measure for free, enabling the
$\widetilde{\omega}_0$ correction of \S\ref{subsec:omega-zero} in
\texttt{RASC-MeetingNotes.tex}.

%%%%%%%%%%%%%%%%%%%%%%%%%%%%%%%%%%%%%%%%%%%%%%%%%%%%%%%%%%%%%%%%%%%%%%%%%%%%
\section*{References}

\begin{thebibliography}{9}

\bibitem{Bertsekas99}
D.~P. Bertsekas.
\newblock {\em Nonlinear Programming}, 2nd ed.
\newblock Athena Scientific, Belmont, MA, 1999.

\bibitem{BecTeb09}
A.~Beck and M.~Teboulle.
\newblock A fast iterative shrinkage-thresholding algorithm for linear inverse
  problems.
\newblock {\em SIAM Journal on Imaging Sciences}, 2(1):183--202, 2009.
\newblock \href{https://doi.org/10.1137/080716542}{doi:10.1137/080716542}.

\bibitem{Bec17}
A.~Beck.
\newblock {\em First-Order Methods in Optimization}.
\newblock SIAM, Philadelphia, PA, 2017.
\newblock \href{https://doi.org/10.1137/1.9781611974997}%
  {doi:10.1137/1.9781611974997}.

\bibitem{BoydADMM11}
S.~Boyd, N.~Parikh, E.~Chu, B.~Peleato, and J.~Eckstein.
\newblock Distributed optimization and statistical learning via the alternating
  direction method of multipliers.
\newblock {\em Foundations and Trends in Machine Learning}, 3(1):1--122, 2011.
\newblock \href{https://doi.org/10.1561/2200000016}{doi:10.1561/2200000016}.

\bibitem{Cha04}
A.~Chambolle.
\newblock An algorithm for total variation minimization and applications.
\newblock {\em Journal of Mathematical Imaging and Vision}, 20(1--2):89--97,
  2004.
\newblock \href{https://doi.org/10.1023/B:JMIV.0000011325.36760.1e}{doi:10.1023/B:JMIV.0000011325.36760.1e}.

\bibitem{ChaPoc11}
A.~Chambolle and T.~Pock.
\newblock A first-order primal-dual algorithm for convex problems with
  applications to imaging.
\newblock {\em Journal of Mathematical Imaging and Vision}, 40(1):120--145,
  2011.
\newblock \href{https://doi.org/10.1007/s10851-010-0251-1}{doi:10.1007/s10851-010-0251-1}.

\bibitem{FleybLey02}
R.~Fletcher and S.~Leyffer.
\newblock Nonlinear programming without a penalty function.
\newblock {\em Mathematical Programming}, 91(2):239--269, 2002.
\newblock \href{https://doi.org/10.1007/s101070100244}{doi:10.1007/s101070100244}.

\bibitem{ParBoy14}
N.~Parikh and S.~Boyd.
\newblock Proximal algorithms.
\newblock {\em Foundations and Trends in Optimization}, 1(3):127--239, 2014.
\newblock \href{https://doi.org/10.1561/2400000003}{doi:10.1561/2400000003}.

\bibitem{RyuBoy16}
E.~K. Ryu and S.~Boyd.
\newblock A primer on monotone operator methods.
\newblock {\em Applied and Computational Mathematics}, 15(1):3--43, 2016.

\end{thebibliography}

\end{document}
